\chapter{Homology}

\begin{center}\itshape
A triangle's three directed edges form a loop. If the triangle is filled in, that loop is the boundary of a face; if only the edges remain, there is nothing inside the space for it to bound. Homology records that difference. It also explains why the boundary of a cube has $8-12+6=2$, while a triangulated torus gives zero. The face counts can change, but their alternating sum survives.
\end{center}

Euler described his polyhedron formula in a letter to Goldbach dated 14 November 1750. Its simplicity concealed a question: which solids should count? A tunnel changes the answer, and a cavity introduces another boundary component. Lhuilier's examples in 1813 made such distinctions hard to ignore. Lakatos later used this history in \emph{Proofs and Refutations} to explore how counterexamples can improve both a theorem and its definitions.

Homology gives a precise way to count cycles that do not bound. Its ingredients are points, paths, triangles and their higher-dimensional counterparts, added with integer coefficients. A closed combination is a cycle. We identify two cycles when their difference bounds a combination one dimension higher. This keeps more information than a count: a cycle may fail to bound while twice that cycle does bound. The resulting element has order two.

Emmy Noether helped make this group viewpoint central in the 1920s. Poincar\'e's Betti numbers and torsion coefficients then became invariants of groups rather than separate lists of numbers. Hopf's 1928 algebraic proof of the Euler--Poincar\'e formula and Alexandroff and Hopf's 1935 \emph{Topologie} reflect that change. The sources at the end give the historical accounts.

We start with finite boundary matrices, where the calculations are visible. Singular chains then let us define the same groups without choosing a triangulation. Subdivision will connect the two descriptions and let us compute by cutting spaces into overlapping pieces.

\section{Chain complexes}

\begin{definition}
A \emph{chain complex} is a sequence of abelian groups $C_k$ and homomorphisms $\partial:C_k\to C_{k-1}$ with $\partial^2=0$; unless an augmentation is specified, $C_k=0$ for $k<0$. Its \emph{cycles} are $Z_k=\ker\partial$, its \emph{boundaries} are $B_k=\partial C_{k+1}$, and its \emph{homology} is
\[
H_k(C)=Z_k/B_k.
\]
A \emph{chain map} $f:C\to C'$ commutes with $\partial$, and therefore induces $f_*:H_k(C)\to H_k(C')$. Maps $f,g$ are \emph{chain homotopic} if maps $P:C_k\to C'_{k+1}$ satisfy $\partial P+P\partial=g-f$.
\end{definition}

The equation $\partial^2=0$ says that every boundary is a cycle, so the quotient makes sense. A chain homotopy makes $f_*=g_*$: for a cycle $z$, the difference $g(z)-f(z)=\partial P(z)$ is a boundary. For example, in a triangular face the three edge endpoints cancel in pairs. The same cancellation in every dimension is the algebra behind our constructions.

A sequence is \emph{exact} when each map's image is the next map's kernel. A \emph{short exact sequence of chain complexes}
$0\to A\xrightarrow{i}B\xrightarrow{j}C\to0$ consists of chain maps and is exact in each degree. It lets a cycle in the quotient $C$ reveal a cycle one dimension lower in $A$.

\begin{lemma}[the zigzag lemma]\label{lem:zigzag}
A short exact sequence of chain complexes $0\to A\xrightarrow{i}B\xrightarrow{j}C\to0$ induces the long exact sequence
\[
\cdots\to H_k(A)\xrightarrow{i_*}H_k(B)\xrightarrow{j_*}H_k(C)
\xrightarrow{\delta}H_{k-1}(A)\to\cdots\to H_0(C)\to0.
\]
For a cycle $c=j(b)$, the connecting map is $\delta[c]=[a]$, where $i(a)=\partial b$. A commutative diagram of short exact sequences induces a commutative diagram of these long sequences.
\end{lemma}

\begin{proof}
Lift the cycle $c$ to $b$. Since $j\partial b=\partial c=0$, exactness gives a unique $a$ with $i(a)=\partial b$. Also $i\partial a=\partial^2b=0$, so $a$ is a cycle. Changing $b$ by $i(a')$ changes $a$ by $\partial a'$. Changing $c$ by a boundary $\partial c'$ permits a change of $b$ by $\partial b'$ with $j(b')=c'$, which leaves $\partial b$ unchanged. These facts prove that $\delta$ is well defined and additive.

At $H_k(B)$, $j_*i_*=0$. If $b$ is a cycle and $j(b)=\partial c'$, lift $c'$ to $b'$; then $b-\partial b'=i(a)$ and injectivity of $i$ makes $a$ a cycle. Thus $[b]=i_*[a]$.

At $H_k(C)$, a projected cycle has zero connecting class. Conversely, if $\delta[c]=0$, write $a=\partial a'$; the lift $b-i(a')$ is a cycle projecting to $c$.

At $H_{k-1}(A)$, $i_*\delta[c]=[\partial b]=0$. Conversely, if $i(a)=\partial b$, then $c=j(b)$ is a cycle and $\delta[c]=[a]$. These three verifications establish exactness at every position, including the end because the complexes vanish in negative degrees. Every choice in the construction can be pushed through a commutative chain diagram, giving naturality.
\end{proof}

\begin{lemma}[the five lemma]\label{lem:five}
In a commutative diagram of abelian groups with exact rows,
\[
\begin{array}{ccccccccc}
A_1&\to&A_2&\to&A_3&\to&A_4&\to&A_5\\
\downarrow\alpha_1&&\downarrow\alpha_2&&\downarrow\alpha_3&&\downarrow\alpha_4&&\downarrow\alpha_5\\
B_1&\to&B_2&\to&B_3&\to&B_4&\to&B_5,
\end{array}
\]
if $\alpha_1,\alpha_2,\alpha_4,\alpha_5$ are isomorphisms, then $\alpha_3$ is an isomorphism too.
\end{lemma}

\begin{proof}
Call the horizontal maps $f_i:A_i\to A_{i+1}$ and $g_i:B_i\to B_{i+1}$. If $\alpha_3(x)=0$, injectivity of $\alpha_4$ gives $f_3(x)=0$, so $x=f_2(y)$. Now $g_2\alpha_2(y)=0$, hence $\alpha_2(y)=g_1(z')$. Write $z'=\alpha_1(z)$; commutativity and injectivity of $\alpha_2$ give $y=f_1(z)$, so $x=0$.

For surjectivity, take $w\in B_3$ and write $g_3(w)=\alpha_4(u)$. Since $\alpha_5f_4(u)=g_4g_3(w)=0$, injectivity of $\alpha_5$ and exactness give $u=f_3(x)$. Then $w-\alpha_3(x)\in\ker g_3=\operatorname{im}g_2$. Surjectivity of $\alpha_2$ writes this difference as $g_2\alpha_2(y)=\alpha_3f_2(y)$. Thus $w=\alpha_3(x+f_2(y))$.
\end{proof}

\section{Simplicial homology}

\begin{definition}
A finite ordered \emph{simplicial complex} $K$ consists of a finite totally ordered vertex set and a family of nonempty vertex subsets, called \emph{simplices}, containing every singleton and every nonempty subset of each member. A simplex $[v_0,\ldots,v_k]$ has $v_0<\cdots<v_k$ and dimension $k$. Realize each vertex as a standard basis vector and each simplex as their convex hull; their union $|K|$ is the \emph{realization}. A \emph{subcomplex} is a subfamily satisfying the same closure rule on its own vertices. A \emph{triangulation} of $X$ is a homeomorphism $|K|\to X$.
\end{definition}

The empty vertex set gives the empty complex and empty realization. For $R=\Z,\Q$ or $\Z/2$, let $C_k^\Delta(K;R)$ be the free $R$-module on the $k$-simplices. Its boundary for $k\ge1$ is
\[
\partial[v_0,\ldots,v_k]=\sum_{i=0}^k(-1)^i[v_0,\ldots,\widehat{v_i},\ldots,v_k].
\]
The hat means omission. Set $\partial=0$ on $C_0$; there is no empty-simplex term in this unaugmented complex. For instance $\partial[v_0,v_1]=[v_1]-[v_0]$ and $\partial[v_0,v_1,v_2]=[v_1,v_2]-[v_0,v_2]+[v_0,v_1]$.

For $i<j$, deleting $v_j$ and then $v_i$ has sign $(-1)^{i+j}$; deleting $v_i$ first shifts $v_j$'s index and gives $(-1)^{i+j-1}$. Each twice-deleted face cancels. Hence $\partial^2=0$. Define $H_k^\Delta(K;R)$ as the homology of these chains. For a subcomplex $L$, use the quotient chains $C^\Delta(K;R)/C^\Delta(L;R)$ to define relative homology $H_k^\Delta(K,L;R)$.

\begin{figure}[ht]
\centering
\fig{ch06-torus}
\caption{A triangulation of the torus: a $3\times3$ grid of squares, each cut by a diagonal, with opposite sides of the big square glued (the vertex labels show the
gluing). It has $V=9$, $E=27$, $F=18$, and $V-E+F=0$. A coarser grid would not do: with a $2\times2$ grid, two different triangles would have the same three vertices.}
\label{fig:torustriang}
\end{figure}

Integer boundary matrices can be reduced by the Smith normal form of Volume~3, Theorem~7.4. Their nonunit invariant factors detect torsion, which ordinary matrix rank cannot see.

\begin{proposition}\label{prop:smithhomology}
Let $c_k$ count the $k$-simplices and $r_k=\operatorname{rank}\partial_k$ over $\Z$. If $d_1\mid\cdots\mid d_{r_{k+1}}$ are the positive nonzero invariant factors of $\partial_{k+1}$, then
\[
H_k^\Delta(K;\Z)\cong\Z^{b_k}\oplus\bigoplus_{d_i>1}\Z/d_i,
\qquad b_k=c_k-r_k-r_{k+1}.
\]
In particular the Euler--Poincar\'e formula is
\[
\chi(K):=\sum_k(-1)^kc_k=\sum_k(-1)^kb_k.
\]
\end{proposition}

\begin{proof}
Choose Smith bases $f_j$ of $C_{k+1}$ and $e_i$ of $C_k$, so $\partial f_i=d_ie_i$ for $i\le r=r_{k+1}$ and the other $f_j$ map to zero. Because $C_{k-1}$ is torsion-free, $d_i\partial e_i=0$ implies $\partial e_i=0$. Thus
\[
Z_k=\bigoplus_{i\le r}\Z e_i\ \oplus\
\left(Z_k\cap\operatorname{span}_{\Z}\{e_{r+1},\ldots,e_{c_k}\}\right).
\]
The second summand is free, by the subgroup theorem for free abelian groups in Volume~3, and has rank $c_k-r_k-r$. Dividing the first summand by $B_k=\bigoplus_{i\le r}d_i\Z e_i$ gives the cyclic factors. Factors with $d_i=1$ vanish. In the finite alternating sum, each $r_k$ occurs twice with opposite signs, giving the formula for $\chi$.
\end{proof}

The integers $b_k$ are the \emph{Betti numbers}. Over a field, Gaussian elimination gives $\dim_RH_k^\Delta=c_k-r_k-r_{k+1}$, with ranks computed over that field. These ranks may change modulo a prime, but the alternating sum is still $\chi(K)$. The displayed torus has $\chi=9-27+18=0$; a tetrahedron's boundary has $4-6+4=2$. We still need to prove that changing the triangulation leaves these answers unchanged.

\section{Singular homology}

Let $\Delta^k$ be the convex hull of $e_0,\ldots,e_k$ in $\R^{k+1}$. A \emph{singular $k$-simplex} is any map $\sigma:\Delta^k\to X$. It need not be injective: it can fold or collapse. Let $C_k(X;R)$ be the free $R$-module on all such maps, with finite sums as its elements. For $k\ge1$ define
\[
\partial\sigma=\sum_{i=0}^k(-1)^i\sigma|[e_0,\ldots,\widehat{e_i},\ldots,e_k],
\]
and again put $\partial_0=0$. Here restriction to an ordered list of vertices means composition with the affine map from the standard face simplex taking its ordered vertices to that list. The same face cancellation proves $\partial^2=0$. Its homology is \emph{singular homology} $H_k(X;R)$, with $R=\Z$ understood when omitted.

A map $f:X\to Y$ sends $\sigma$ to $f\sigma$, a chain map $f_\#$. The induced maps obey $(gf)_*=g_*f_*$ and $\mathrm{id}_*=\mathrm{id}$. Thus homeomorphic spaces have isomorphic homology without any choice of triangulation.

For $A\subset X$, singular simplices in $A$ give a subcomplex $C(A;R)$. Define $H_k(X,A;R)$ using $C(X;R)/C(A;R)$. A relative cycle is a chain whose boundary lies in $A$. The zigzag lemma gives the natural sequence
\[
\cdots\to H_k(A)\to H_k(X)\to H_k(X,A)\xrightarrow{\delta}H_{k-1}(A)\to\cdots.
\]

To remove the extra copy of $R$ contributed by a nonempty connected space, \emph{augment} its chains: add $C_{-1}=R$ and replace $\partial_0$ by $\varepsilon(\sum a_x[x])=\sum a_x$. Homology of this augmented complex is \emph{reduced homology} $\tilde H_k(X;R)$. It equals $H_k$ for $k>0$. For nonempty $X$, choosing one point splits the augmentation and gives $H_0(X)\cong\tilde H_0(X)\oplus R$, and $\tilde H_{-1}(X)=0$. For the empty space, all unaugmented chains vanish, while $\tilde H_{-1}(\emptyset)=R$ and its other reduced groups vanish. All groups below degree $-1$ vanish. This negative-degree convention will make the sphere calculations work at their endpoints.

\begin{proposition}\label{prop:H0}
The module $H_0(X;R)$ is free on the path components of $X$. A point has zero reduced homology in every degree.
\end{proposition}

\begin{proof}
A singular $0$-simplex is a point. The boundary of a path from $x$ to $y$ is $[y]-[x]$. Sending each point to its path component therefore kills boundaries and maps onto the free module on those components. Its kernel consists of finite point sums with total coefficient zero in each component. Join each of the finitely many points to a chosen point of its component; their differences are path boundaries, so the kernel is precisely $B_0$. This is a direct sum, even with infinitely many components.

For a one-point space there is one singular simplex $\sigma_k$ in each degree. Its boundary is $\sigma_{k-1}$ for positive even $k$ and zero for odd $k$, since the alternating sum of $k+1$ signs is respectively one or zero. The augmentation $C_0\to R$ is the identity. This alternating complex is exact, proving the reduced assertion.
\end{proof}

\begin{theorem}[homotopy invariance]\label{thm:homotopyH}
Homotopic maps $f,g:X\to Y$ induce the same homology maps. Hence a homotopy equivalence induces homology isomorphisms, and a contractible space has zero reduced homology.
\end{theorem}

\begin{proof}
Let $F:X\times I\to Y$ join $f$ to $g$. For $\sigma:\Delta^k\to X$ write $v_i=(e_i,0)$, $w_i=(e_i,1)$ and $\Phi=F\circ(\sigma\times\mathrm{id})$. Define the \emph{prism operator}
\[
P\sigma=\sum_{i=0}^k(-1)^i\Phi|[v_0,\ldots,v_i,w_i,\ldots,w_k].
\]
For $k=0$ this is the path tracing the homotopy of a point. For general $k$, the terms triangulate the prism over $\Delta^k$ with consistent signs. We check the boundary identity rather than assuming it from the picture.

Deleting a bottom vertex $v_j$, $j\le i$, contributes sign $(-1)^{i+j}$. Deleting a top vertex $w_j$, $j\ge i$, contributes $(-1)^{i+j+1}$. The diagonal face obtained by deleting $v_i$ in prism $i$ has sign $+1$ and cancels the same face obtained by deleting $w_{i-1}$ in prism $i-1$, with sign $-1$. At the two ends this leaves $g\sigma-f\sigma$. All other faces lie over a proper face of $\Delta^k$. If its missing index is $j<i$, its prism index becomes $i-1$, and its sign in $P\partial\sigma$ is $(-1)^{j+i-1}$, opposite to the boundary sign. If $j>i$, its prism index stays $i$, giving sign $(-1)^{j+i}$, again opposite. Thus
\[
\partial P+P\partial=g_\#-f_\#.
\]
This chain homotopy proves $f_*=g_*$. Homotopy inverses then induce inverse homology maps. A contraction is a homotopy equivalence to a point, whose reduced groups we have computed.
\end{proof}

The construction respects subspaces: if $F(A\times I)\subset B$, the prism of a chain in $A$ lies in $B$. Thus homotopy invariance also holds for maps of pairs. The pair sequence shows that $H_k(X,A)=0$ when $A$ retracts by deformation from $X$. If $A\subset A'\subset X$ and $A$ is a deformation retract of $A'$, comparing the two pair sequences and using the five lemma gives $H_k(X,A)\cong H_k(X,A')$.

\section{Small chains and Mayer--Vietoris}

A large simplex can cross both pieces of a cover. Cutting it into smaller simplices will let us separate a cycle into a chain in each piece. The cuts must preserve its homology class, including the cancellation along shared faces.

Let $\mathcal U=\{U_j\}$ be subsets whose interiors cover $X$. Write $C^{\mathcal U}(X)$ for chains all of whose simplices have image in some $U_j$. The face of such a simplex stays in the same set, so these chains form a subcomplex.

\begin{theorem}[small chains]\label{thm:smallchains}
The inclusion $C^{\mathcal U}(X)\hookrightarrow C(X)$ induces homology isomorphisms.
\end{theorem}

\begin{proof}
First work with affine singular simplices $\lambda=[w_0,\ldots,w_k]$ in a convex subset of Euclidean space; repeated vertices are allowed. For $b$ in that set, the cone $b\lambda=[b,w_0,\ldots,w_k]$ obeys
\[
\partial(b\lambda)=\lambda-b\partial\lambda\quad(k\ge1),\qquad
\partial(b\lambda)=\lambda-\varepsilon(\lambda)[b]\quad(k=0).
\]
These follow by deleting the first vertex and then the remaining vertices with their signs. Let $b_\lambda$ be the barycentre. Define subdivision $S$ and a homotopy $T$ inductively:
\[
S\lambda=\lambda\ (k=0),\quad S\lambda=b_\lambda(S\partial\lambda)\ (k>0),
\qquad
T\lambda=0\ (k=0),\quad T\lambda=b_\lambda(\lambda-T\partial\lambda)\ (k>0).
\]
The cone identity and induction give $\partial S=S\partial$. For $k=1$ the augmentation terms in $S\partial\lambda$ cancel, as the coefficients of a boundary sum to zero. In higher degrees $\partial S\partial\lambda=0$ gives the same conclusion. For $T$, the induction identity on the boundary gives
\[
\partial T\lambda
=\lambda-T\partial\lambda-b_\lambda(\partial\lambda-\partial T\partial\lambda)
=\lambda-T\partial\lambda-b_\lambda S\partial\lambda
=\lambda-T\partial\lambda-S\lambda.
\]
Thus $\partial T+T\partial=\mathrm{id}-S$. It also holds in degree zero because $S=\mathrm{id}$ there.

The vertices of each subdivided simplex are barycentres of a nested sequence of faces. If $F\subset G$ and $G$ has dimension $j\le k$, then
\[
|b_G-b_F|\le\max_{w\text{ vertex of }G}|b_G-w|
\le\frac{j}{j+1}\operatorname{diam}G
\le\frac{k}{k+1}\operatorname{diam}\lambda.
\]
The middle bound follows from $b_G-w=(j+1)^{-1}\sum_{w'\ne w}(w'-w)$. Hence subdivision in positive dimension shrinks diameters by at most $k/(k+1)$; repeated subdivision makes them tend to zero.

For arbitrary $\sigma:\Delta^k\to X$, perform these constructions on the identity simplex of $\Delta^k$ and compose all terms with $\sigma$. Affine maps preserve barycentres, so face restrictions commute with the construction. We obtain the same identities $\partial S=S\partial$ and $\partial T+T\partial=\mathrm{id}-S$ on singular chains. Both operators keep every term inside its original simplex image. With
\[
D_m=\sum_{i=0}^{m-1}TS^i
\]
we have $\partial D_m+D_m\partial=\mathrm{id}-S^m$, also for $m=0$ with an empty sum.

Given a finite chain, choose $m$ large enough geometrically for all its simplices: the sets $\sigma^{-1}(\operatorname{int}U_j)$ cover each compact domain simplex, so a Lebesgue number and the diameter estimate put every subdivided piece inside some $U_j$. Choose this depth before any cancellation of terms. Zero-dimensional simplices are already small.

If $z$ is a cycle, such an $m$ makes $S^mz$ small and $z-S^mz=\partial D_mz$, proving surjectivity on homology. If a small cycle $z$ bounds $b$ in $C(X)$, choose $m$ making $S^mb$ small. Since $D_m$ preserves supports, $D_mz$ is small too, and
\[
z=S^mz+\partial D_mz=\partial(S^mb+D_mz)
\]
is a boundary in the small complex. This proves injectivity.
\end{proof}

\begin{theorem}[Mayer--Vietoris]\label{thm:mv}
If $X=\operatorname{int}A\cup\operatorname{int}B$, there is a natural long exact sequence
\[
\cdots\to H_k(A\cap B)\xrightarrow{(i_{A*},i_{B*})}H_k(A)\oplus H_k(B)
\xrightarrow{j_{A*}-j_{B*}}H_k(X)\xrightarrow{\delta}H_{k-1}(A\cap B)\to\cdots.
\]
Except for $\delta$, the maps come from inclusions. If a cycle is written $z=x+y$, with $x$ in $A$ and $y$ in $B$, then $\delta[z]=[\partial x]$. When $A\cap B\ne\emptyset$, the same sequence holds in reduced homology.
\end{theorem}

\begin{proof}
Use the short exact sequence
\[
0\to C(A\cap B)\xrightarrow{c\mapsto(c,c)}C(A)\oplus C(B)
\xrightarrow{(a,b)\mapsto a-b}C^{\{A,B\}}(X)\to0.
\]
A singular simplex belongs to both chain bases exactly when its image lies in $A\cap B$, giving exactness at the middle. The last map is onto by the definition of small chains. Apply the zigzag lemma and the small-chains isomorphism. For the connecting formula, $z=x+y$ lifts to $(x,-y)$, whose boundary is $(\partial x,-\partial y)=(\partial x,\partial x)$. For the reduced version, augment by the exact degree $-1$ sequence $0\to R\to R\oplus R\to R\to0$ with these same diagonal and difference maps. Naturality follows from the chain constructions.
\end{proof}

\begin{corollary}[excision]\label{cor:excision}
If $Z\subset A\subset X$ and $\overline Z\subset\operatorname{int}A$, inclusion induces
\[
H_k(X\setminus Z,A\setminus Z)\cong H_k(X,A).
\]
\end{corollary}

\begin{proof}
Put $B=X\setminus Z$. The interior of $B$ contains $X\setminus\overline Z$, so the interiors of $A,B$ cover $X$. Modulo chains in $A$, every small chain for this cover is a chain in $B$, giving the chain isomorphism
\[
C^{\{A,B\}}(X)/C(A)\cong C(B)/C(A\cap B).
\]
Compare the long exact sequences for $C(A)\subset C^{\{A,B\}}(X)$ and $C(A)\subset C(X)$. Their maps on $H(A)$ are identities and those on the middle complexes are small-chain isomorphisms. The five lemma gives the required quotient homology isomorphism, induced by inclusion.
\end{proof}

\section{Spheres, degree and the agreement of the two theories}

\begin{theorem}\label{thm:spheres}
For $n\ge0$, $\tilde H_k(S^n;R)$ is $R$ when $k=n$ and zero otherwise. For $n\ge1$, let $U_\pm=\{x\in S^n:\pm x_{n+1}>-1/2\}$, with equator $S^{n-1}=\{x_{n+1}=0\}$. The map
\[
\delta:\tilde H_k(S^n)\longrightarrow\tilde H_{k-1}(U_+\cap U_-)
\cong\tilde H_{k-1}(S^{n-1})
\]
is an isomorphism, natural for maps preserving both caps and the equator.
\end{theorem}

\begin{proof}
For $S^0=\{-1,1\}$, reduced $H_0$ is generated by $[1]-[-1]$ and is $R$; all other reduced groups vanish by the point calculation and the component decomposition. For $n\ge1$, each cap contracts to its pole, for example by normalizing straight segments to that pole. The overlap retracts onto the equator by shrinking the last coordinate and normalizing. Its horizontal part never vanishes in this band. Reduced Mayer--Vietoris, whose cap terms vanish, gives the displayed isomorphism. Induct on $n$. Its naturality comes from Mayer--Vietoris and the equatorial inclusion, which is a homotopy equivalence onto the band.
\end{proof}

For $n\ge0$, a map $f:S^n\to S^n$ multiplies $\tilde H_n(S^n;\Z)=\Z$ by an integer, its \emph{degree}. For $n>0$ this is ordinary top homology. Homotopy invariance gives homotopy invariance of degree, and composition gives $\deg(fg)=\deg f\deg g$, with $\deg\mathrm{id}=1$. A map missing a point has degree zero, because it factors through the contractible punctured sphere; for $S^0$ such a map is constant, so the same assertion follows directly. We will identify the circle's degree with winding below.

\begin{proposition}\label{prop:reflection}
A reflection of $S^n$, for example $r(x_1,\ldots,x_{n+1})=(-x_1,x_2,\ldots,x_{n+1})$, has degree $-1$. The antipodal map has degree $(-1)^{n+1}$.
\end{proposition}

\begin{proof}
For $n\ge1$, $r$ preserves the caps and equator in the preceding theorem. Its action commutes with the connecting isomorphism, reducing its degree to that of the same reflection on $S^{n-1}$. Eventually we reach $S^0$, where it exchanges the points and negates $[1]-[-1]$. Reflections in the other coordinate hyperplanes are conjugate to this one by coordinate permutations. Conjugation does not change degree, since a homeomorphism has degree $1$ or $-1$. Multiplying the $n+1$ coordinate reflections gives the antipodal map and the asserted degree.
\end{proof}

\begin{corollary}\label{cor:brouwerH}
For $n\ge0$, the closed ball $\bar B^{n+1}$ does not retract onto $S^n$, and every map from that ball to itself has a fixed point. If $m\ne n$ and $m,n\ge1$, no nonempty open subset of $\R^m$ is homeomorphic to an open subset of $\R^n$.
\end{corollary}

\begin{proof}
A retraction $r$ and inclusion $i$ would give $r_*i_*=\mathrm{id}$ on $\tilde H_n(S^n)=\Z$, factoring through the zero reduced homology of the ball. For the fixed-point assertion, if $f$ had no fixed point, send $x$ to the boundary point on the ray from $f(x)$ through $x$. The ray has a unique forward boundary intersection, continuously depending on its two distinct points, and it fixes boundary points $x$. This would be a retraction, a contradiction. The quadratic ray formula was also used in Theorem~\ref{thm:brouwern}; no differentiability is needed.

For $x$ in open $U\subset\R^m$, excising the complement of a small open ball about $x$ identifies $H_k(U,U\setminus\{x\})$ with $H_k(\R^m,\R^m\setminus\{x\})$. The pair sequence and the radial retraction of punctured Euclidean space onto $S^{m-1}$ give $\Z$ for $k=m$ and zero otherwise. At $k=0$, the map on $H_0$ from the nonempty punctured space to $\R^m$ is onto, so the relative group is zero as required. A homeomorphism preserves these local relative groups, forcing $m=n$. The omitted dimension-zero comparison is immediate: $\R^0$ is a single point, while a positive-dimensional nonempty open set has no isolated points. The fixed-point theorem for its zero-dimensional ball is also immediate.
\end{proof}

Now associate to an ordered simplex $[v_0,\ldots,v_k]$ its affine singular simplex in $|K|$. This defines a chain map $\varphi:C^\Delta(K;R)\to C(|K|;R)$, compatible with subcomplex inclusions.

\begin{theorem}\label{thm:simplicialsingular}
For every finite simplicial complex $K$, $\varphi_*:H_k^\Delta(K;R)\to H_k(|K|;R)$ is an isomorphism in every degree. The same holds for pairs $(K,L)$.
\end{theorem}

\begin{proof}
Induct on the number of simplices, starting with the empty complex, whose unaugmented chains vanish on both sides. Choose a maximal simplex $\sigma$ of dimension $d$, and let $L$ consist of all other simplices. It is a subcomplex because $\sigma$ is not a proper face of another simplex. Write $\bar\sigma$ for the complex of all its faces and $\partial\bar\sigma$ for its proper faces.

The comparison holds for the single simplex $\bar\sigma$ without induction. Its singular realization is convex. On simplicial augmented chains, coning to the first vertex, with a repeated first vertex interpreted as zero, gives $\partial h+h\partial=\mathrm{id}$ by the face formula. Equivalently its unaugmented homology is $R$ in degree zero and zero above, and $\varphi_*$ sends a vertex generator to the corresponding point generator. The boundary complex has fewer simplices than $K$, so induction and the pair sequences, with the five lemma, give the comparison for $(\bar\sigma,\partial\bar\sigma)$.

On simplicial relative chains, inclusion from this pair to $(K,L)$ is an isomorphism: only the generator $\sigma$ in degree $d$ remains. We verify the singular relative inclusion too. For $d=0$, $\sigma$ is an isolated point, and the singular chains split as the direct sum for that point and $|L|$, proving the claim directly. For $d>0$, let $b$ be the barycentre and $W=|K|\setminus\{b\}$. Radial projection from $b$ retracts $W$ onto $|L|$, fixing $|L|$; within $|\sigma|$ it retracts $|\sigma|\setminus\{b\}$ onto $|\partial\bar\sigma|$. The maps from the original pairs to these punctured pairs therefore induce isomorphisms.

Excision removes the closed set $|L|$ from $(|K|,W)$ and the closed boundary $|\partial\bar\sigma|$ from $(|\sigma|,|\sigma|\setminus\{b\})$. Both removed sets lie inside the relevant open punctured subspaces. The results are the same pair
\[
(\sigma^\circ,\sigma^\circ\setminus\{b\}).
\]
Here $\sigma^\circ$ is open in $|K|$ because $\sigma$ is maximal in a finite complex. Every map is induced by inclusion, so these identifications show that the singular inclusion of pairs is an isomorphism. In the comparison square, the other three maps are now isomorphisms, hence so is the map for $(K,L)$. The pair sequences and the induction hypothesis for $L$ give the result for $K$ by the five lemma. Finally, for any subcomplex pair, the two absolute comparisons and its pair sequences give the relative comparison.
\end{proof}

\begin{corollary}\label{cor:euler}
The Betti numbers and Euler characteristic of a finite simplicial complex depend only on the homotopy type of its realization. Thus all finite triangulations of a space have the same alternating simplex count. In particular, the boundary of every bounded full-dimensional convex polyhedron in $\R^3$ satisfies $V-E+F=2$.
\end{corollary}

\begin{proof}
The comparison theorem identifies the Betti numbers with ranks of singular homology, which is homotopy invariant. Their alternating sum gives $\chi$. For the polyhedron, triangulate its polygonal faces by diagonals: each diagonal adds one edge and one face, preserving $V-E+F$. Radial projection from an interior point identifies the boundary with $S^2$, whose Betti numbers are $1,0,1$. Hence the alternating count is two. Boundedness and nonempty three-dimensional interior ensure that this boundary really is a sphere; lower-dimensional or unbounded convex polyhedra need not have that boundary.
\end{proof}

\begin{example}\label{ex:surfacesH}
\emph{Turning loops into homology classes.} A loop is a singular $1$-cycle. For based loops $\gamma,\eta$, compose $\gamma\eta$ with the affine map $\Delta^2\to I$ taking its vertices to $0,1/2,1$. Its boundary is $\eta-\gamma\eta+\gamma$. Thus concatenation becomes addition in $H_1$. A based homotopy gives a homology by the prism operator; its constant side paths cancel. We obtain a homomorphism $h:\pi_1(X,x_0)\to H_1(X)$.

On $S^1$, cover by two open arcs with intersection two arcs. Mayer--Vietoris gives $H_1(S^1)\cong\tilde H_0(\text{two points})=\Z$. Subdividing a positive loop into its arc pieces sends it under $\delta$ to the difference of the two overlap components, a generator up to sign. Hence $h$ identifies the winding generator of $\pi_1(S^1)$ with a generator of $H_1(S^1)$. By naturality, a circle map's action on $H_1$ is multiplication by its winding number, so this is its degree.

\emph{Surface groups become abelian.} Use the bouquet open covers from Chapter~5 with Mayer--Vietoris. Their contractible star intersections show that a bouquet of $m$ circles has $H_1=\Z^m$, with its loops as generators, and all higher homology zero. For $\Sigma_g$, remove the polygon disc's centre to obtain $U$, retracting onto the bouquet of $2g$ circles, and take its open interior as $V$. The intersection retracts onto its boundary circle. The attaching word has homology class
\[
a_1+b_1-a_1-b_1+\cdots+a_g+b_g-a_g-b_g=0.
\]
Reduced Mayer--Vietoris therefore contains
\[
0\to H_2(\Sigma_g)\to\Z\xrightarrow{0}\Z^{2g}
\to H_1(\Sigma_g)\to0.
\]
It follows that $H_2=\Z$, $H_1=\Z^{2g}$, $H_0=\Z$, and all higher groups vanish. Thus $\chi(\Sigma_g)=2-2g$. For $g=0$ the sphere computation gives the same formula directly. This is also the index sum found in Chapter~4 for a smooth surface with $g$ handles.
\end{example}

\begin{remark}
We quote two surface theorems here. Every compact surface is triangulable (Rad\'o, 1925), and every closed connected orientable surface is homeomorphic to exactly one $\Sigma_g$ (the classification associated with Dehn and Heegaard, 1907, and Brahana, 1921). We do not prove either theorem. They let us identify the Euler characteristic from triangulation with the smooth index sum of Chapter~4: on each connected closed orientable surface both are $2-2g$. For a disconnected compact surface, its chain complex and index sum split over its finitely many components, so their values add. Hopf's 1927 paper gives a direct vector-field proof using vertices, edges and faces.
\end{remark}

\begin{figure}[ht]
\centering
\fig{ch06-mv}
\caption{Mayer--Vietoris on the sphere: two open caps $U_+$ and $U_-$, each contractible, overlap in a band that retracts onto the equator. A cycle on $S^n$ is the sum of a chain in $U_+$
and a chain in $U_-$; its image under $\delta$ is their common boundary, a cycle in the band, one dimension lower. So every sphere carries exactly one hole, in its top dimension.}
\label{fig:mvsphere}
\end{figure}

\begin{figure}[ht]
\centering
\fig{v9-ch06-filled-cycle}
\caption{The same oriented three-edge cycle in two spaces. On the left it generates $H_1\cong\Z$; on the right it is the boundary $\partial[0,1,2]=[1,2]-[0,2]+[0,1]$ of the filled triangle, so its class is zero. Homology asks what bounds inside the space.}
\label{fig:v9-filled-cycle}
\end{figure}

\begin{figure}[ht]
\centering
\fig{v9-ch06-local-stars}
\caption{Remove the central vertex of a small star. One, two and three arms leave one, two and three components. The local relative $H_1$ groups are consequently $0$, $\Z$ and $\Z^2$. This distinguishes an interval endpoint, an interior point and a branch point; see Exercise~\ref{ex:v9-local-homology}.}
\label{fig:v9-local-stars}
\end{figure}

\section{Invariance of domain and Borsuk--Ulam}

\begin{proposition}\label{prop:jordanbrouwer}
Let $n\ge1$.
\begin{enumerate}
\item[(a)] If $h:\bar B^k\to S^n$ is an embedding, then $\tilde H_i(S^n\setminus h(\bar B^k))=0$ in every degree.
\item[(b)] If $h:S^k\to S^n$ is an embedding with $0\le k<n$, then $\tilde H_i(S^n\setminus h(S^k))$ is $\Z$ for $i=n-k-1$ and zero otherwise.
\end{enumerate}
In particular an embedded $S^{n-1}$ has a complement with exactly two path components, whereas an embedded closed ball has path-connected complement. An embedding means a homeomorphism onto its image; for these compact domains, a continuous injection into $S^n$ is an embedding.
\end{proposition}

\begin{proof}
(a) Replace the ball by the homeomorphic cube $I^k$. Its complement is nonempty: otherwise the embedding would identify a contractible cube with the noncontractible sphere, contrary to their homology. For $k=0$ the complement of a point is $\R^n$, so the assertion holds.

Induct on $k$. Suppose a nonzero reduced class $\alpha$ exists in the complement of $h(I^k)$. Cut the last coordinate interval in half. The complements of the two image half-cubes are open sets $A,B$ with
\[
A\cap B=S^n\setminus h(I^k),\qquad
A\cup B=S^n\setminus h(I^{k-1}\times\{1/2\}).
\]
Induction makes the reduced homology of the union zero. Mayer--Vietoris therefore makes $\tilde H_i(A\cap B)\to\tilde H_i(A)\oplus\tilde H_i(B)$ an isomorphism, so $\alpha$ remains nonzero in at least one half-cube complement. Repeating produces nested closed intervals $J_m$ of lengths $2^{-m}$, with intersection $\{p\}$, such that the class stays nonzero in each $S^n\setminus h(I^{k-1}\times J_m)$.

These increasing open complements have union $S^n\setminus h(I^{k-1}\times\{p\})$. The latter has zero reduced homology by induction, so a cycle representing $\alpha$ bounds a finite chain there. Its compact simplex images lie in one of the increasing complements: a finite subcover by these nested open sets suffices. It would then bound there already, a contradiction. In degree zero the cycle has augmentation zero and the same bounding argument applies; negative reduced degrees vanish because the complements are nonempty.

(b) For $k=0$, remove two points from $S^n$. The complement is $\R^n$ minus a point and retracts onto $S^{n-1}$. For $k\ge1$, divide $S^k$ into closed hemispheres $D_+,D_-$ meeting in $S^{k-1}$. The open sets $A=S^n\setminus h(D_+)$ and $B=S^n\setminus h(D_-)$ have zero reduced homology by (a). Their union is the complement of the equator image, and their intersection is the complement of $h(S^k)$. Mayer--Vietoris gives
\[
\tilde H_i(S^n\setminus h(S^k))
\cong\tilde H_{i+1}(S^n\setminus h(S^{k-1})).
\]
Induction proves the formula. Finally $\tilde H_0$ is free with one fewer generator than the number of path components when that number is finite; zero reduced $H_0$ for a nonempty space means one component, and $\tilde H_0=\Z$ means two. All complements here are open in $S^n$, hence locally path-connected, so these path components are open as well.
\end{proof}

\begin{theorem}[invariance of domain; Brouwer, 1912]\label{thm:invariancedomain}
If $U\subset\R^n$ is open and $h:U\to\R^n$ is continuous and injective, then $h(U)$ is open and $h$ is a homeomorphism onto its image.
\end{theorem}

\begin{proof}
For $n=0$ the assertion concerns subsets of a single point and is immediate; the empty domain in any dimension is also immediate. Let $n\ge1$ and $x\in U$. Choose a closed ball $D$ about $x$ inside $U$, with boundary $S$. Regard the target $\R^n$ as a punctured $S^n$. The compact restriction to $D$ is an embedding. The preceding proposition makes $S^n\setminus h(S)$ have two open path components. It is the disjoint union of the nonempty path-connected sets
\[
S^n\setminus h(D)\quad\hbox{and}\quad h(D\setminus S).
\]
Each lies in one component. Since they cover both components, they must be precisely those two components. Thus $h(D\setminus S)$ is open in $S^n$, lies in $h(U)$ and contains $h(x)$. The image $h(U)$ is open. Applying this same argument to every open subset of $U$ shows that $h$ is an open map, so its inverse on the image is continuous.
\end{proof}

\begin{beauty}
The hypothesis says only that $h$ is continuous and one-to-one. It gives no derivative and no local inverse formula. Nevertheless, a sphere around $x$ still separates the target after deformation, and the image of the ball's interior must occupy one of the two components. Homology of the complement supplies the neighbourhood that calculus would have obtained from an invertible derivative. Volume~7, Chapter~3 asks whether $\R^2$ and $\R^3$ are homeomorphic. Deleting one point leaves both complements path-connected, and the comparison is left to this volume. Corollary~\ref{cor:brouwerH} separates the dimensions by local homology. The openness just proved separates them by a second route. If a homeomorphism carried $\R^3$ onto $\R^2$, inclusion of that coordinate plane in $\R^3$ would turn it into a continuous injection of the open set $\R^3$ into $\R^3$. The image would be a nonempty open subset of $\R^3$ lying in the plane. Every open ball in $\R^3$ meets the complement of the plane, so no such open set exists.
\end{beauty}

For the antipodal problem, coefficients modulo two let us add a simplex's two lifts without choosing a sign. All chain and gluing arguments above apply with these coefficients.

\begin{lemma}\label{lem:rpn}
For $n\ge0$, $H_k(\R P^n;\Z/2)=0$ when $k>n$.
\end{lemma}

\begin{proof}
The case $n=0$ is a point. For $n\ge1$, remove $p=[0:\cdots:0:1]$ to form $U$. The homotopy
\[
[x_0:\cdots:x_n]\longmapsto[x_0:\cdots:x_{n-1}:(1-t)x_n]
\]
retracts $U$ onto $\R P^{n-1}$; its first $n$ coordinates cannot all vanish in $U$. Let $V=\{[x]:|x_n|>|x|/2\}$. In the chart $x_n=1$ it is the open ball $\sum_{j<n}x_j^2<3$, and $U\cap V=V\setminus\{p\}$ retracts onto $S^{n-1}$. For $k>n$, the outer terms in
\[
H_k(U)\oplus H_k(V)\to H_k(\R P^n)\to H_{k-1}(U\cap V)
\]
vanish by induction, contractibility and the sphere calculation. Exactness gives the claim.
\end{proof}

\begin{theorem}[Borsuk--Ulam]\label{thm:borsukulam}
For $n\ge1$, no continuous map $g:S^n\to S^{n-1}$ is odd, meaning $g(-x)=-g(x)$. Equivalently, every continuous $f:S^n\to\R^n$ has $f(x)=f(-x)$ for some $x$.
\end{theorem}

\begin{proof}
For $n=1$, a map from connected $S^1$ to the two-point space $S^0$ is constant and cannot be odd. Take $n\ge2$ and use $\Z/2$ coefficients for the remainder of the homology argument.

Every singular simplex of $\R P^n$ has two lifts through $q:S^n\to\R P^n$, since its domain is convex and simply connected. Define the \emph{transfer} $\tau\sigma$ to be the sum of those two lifts. Faces lift in the same way, so $\tau$ is a chain map. We have a short exact sequence
\[
0\to C(\R P^n)\xrightarrow{\tau}C(S^n)\xrightarrow{q_\#}C(\R P^n)\to0.
\]
Indeed $q_\#$ is onto, $q_\#\tau=2=0$, and a chain maps to zero exactly when the coefficients of its two lifts over every simplex agree. Such chains are precisely transfer images. Distinct downstairs simplices have disjoint pairs of lifts, so $\tau$ is injective.

The zigzag lemma yields
\[
\cdots\to H_k(\R P^n)\xrightarrow{\tau_*}H_k(S^n)
\xrightarrow{q_*}H_k(\R P^n)\xrightarrow{\delta}H_{k-1}(\R P^n)\to\cdots.
\]
The map $q_*$ on $H_0=\Z/2$ is the identity because both spaces are path-connected. Exactness then makes $\tau_*=0$ on $H_0$, so $\delta:H_1(\R P^n)\to H_0(\R P^n)$ is an isomorphism. In the remaining degrees $1<k<n$, the sphere terms on either side vanish, making $\delta$ an isomorphism again. Hence $H_k(\R P^n)=\Z/2$ for $0\le k<n$.

At the top the preceding lemma gives the exact segment
\[
0\to A\xrightarrow{\tau_*}\Z/2\xrightarrow{q_*}A
\xrightarrow{\delta}\Z/2\to0,\qquad A=H_n(\R P^n).
\]
The injection restricts $A$ to zero or $\Z/2$, while the final surjection rules out zero. Thus $A=\Z/2$, the transfer is an isomorphism, $q_*=0$, and $\delta$ is an isomorphism. These specific connecting and transfer maps are isomorphisms; the quotient map at the top is not.

If an odd $g:S^n\to S^{n-1}$ existed, include its target as the equator to obtain an odd self-map $f$ of $S^n$. It induces $\bar f$ on projective space. Oddness gives $f_\#\tau=\tau\bar f_\#$, because it carries each pair of antipodal lifts to the pair over the composed simplex. Naturality makes $\bar f_*$ commute with $\delta$. It is the identity on $H_0$, hence on every $H_k(\R P^n)$ through the connecting isomorphisms. The top transfer then makes $f_*$ an isomorphism on $H_n(S^n)$. But $f$ misses a point, so it factors through a contractible punctured sphere and has zero action there. This contradiction excludes $g$.

Finally, if $f(x)-f(-x)$ never vanished, its normalization would be an odd map to $S^{n-1}$. Conversely an odd map into $S^{n-1}\subset\R^n$ has distinct antipodal values, since its values are nonzero. This proves the equivalence.
\end{proof}

The computation also gives $H_k(\R P^n;\Z/2)=\Z/2$ for $0\le k\le n$ and zero above. For $n=1$ this follows from $\R P^1\cong S^1$, and for $n=0$ from the point calculation. Modulo two, projective space has a nonzero homology class in each degree up to its dimension.

\begin{corollary}[the ham sandwich theorem]\label{cor:hamsandwich}
For $n\ge1$, any $n$ bounded Lebesgue measurable subsets $A_1,\ldots,A_n$ of $\R^n$ can be bisected simultaneously by one affine hyperplane, with equal measure on its two open sides.
\end{corollary}

\begin{proof}
For $u=(u_0,u')\in S^n\subset\R\times\R^n$, put
\[
f_i(u)=\lambda\{x\in A_i:\langle u',x\rangle>u_0\}.
\]
When $u'\ne0$, the boundary hyperplane has measure zero, by a linear change of variables and Fubini. For convergent parameters its half-space indicator therefore converges almost everywhere; dominated convergence, with the integrable bound $1_{A_i}$, proves continuity. At $(1,0)$ the half-space misses all bounded $A_i$ for nearby parameters, and at $(-1,0)$ it contains them, so continuity holds there too.

Borsuk--Ulam gives $f_i(u)=f_i(-u)$ for every $i$. For $u'\ne0$ the two open half-spaces exhaust each $A_i$ up to a null hyperplane, so both measures equal $\lambda(A_i)/2$. If $u'=0$, equality forces every $A_i$ to be null, in which case any affine hyperplane works.
\end{proof}

In three dimensions, one plane can bisect bread, ham and cheese at once, regardless of their bounded measurable shapes. Steinhaus posed the problem in Lw\'ow, and Banach reduced it to the antipodal theorem; the 1938 \emph{Mathesis Polska} note records that exchange. Stone and Tukey published the general-dimensional version in 1942.

\section{What we proved, and what comes next}

The filled triangle from the opening carries a chain whose boundary is its edge loop. Removing the face removes that boundary without removing the cycle. In larger spaces the same quotient of cycles by boundaries gives groups that survive homotopy. Singular chains make them topological invariants; the comparison theorem lets finite triangulations compute them by matrices.

On $\Sigma_g$, the calculation gives one connected component, $2g$ independent one-dimensional classes and one top-dimensional class. Their alternating ranks are $1-2g+1=2-2g$. That explains why triangulating a handle changes the Euler count without making it depend on the chosen faces. Torsion adds information that the alternating rank cannot retain, as the projective-plane exercises show.

Part~II ends with these ways of distinguishing spaces. Part~III asks what happens when a manifold is also a group. Rotations and invertible matrices can be multiplied as well as moved continuously. Chapter~7 uses that extra operation to connect tangent vectors, exponentials and global transformations.

\section{Exercises}

\begin{exercise}[$\star$]\label{ex:kleinrp2}
\emph{Torsion from a boundary word.} (a) Present $\R P^2$ as a disc attached to a circle by $aa$. Compute $H_1,H_2$ with integer coefficients and $H_2$ modulo two. (b) For the Klein bottle's square with boundary word $aba^{-1}b$, prove $H_1=\Z\oplus\Z/2$, $H_2=0$ and $\chi=0$.
\end{exercise}

\begin{exercise}[$\star\star$]\label{ex:hairydegree}
\emph{A continuous hairy-ball proof.} Let $n\ge1$. (a) For a continuous tangent field $X$ without zeros on $S^n$, check that
\[
H(x,t)=\cos(\pi t)x+\sin(\pi t)\frac{X(x)}{|X(x)|}
\]
is a homotopy from the identity to the antipodal map. (b) Deduce that $n$ is odd. (c) Prove that a fixed-point-free map $f:S^n\to S^n$ is homotopic to the antipodal map, while a map with $f(x)\ne-x$ everywhere is homotopic to the identity. Conclude that for even $n$ every self-map has a fixed point or an antipodal value. Discuss $S^0$ separately.
\end{exercise}

\begin{exercise}[$\star$]\label{ex:lhuilier}
\emph{Counting a frame's boundary.} (a) A square picture frame has a torus as its boundary. Verify $V-E+F=0$ in Figure~\ref{fig:torustriang} and explain why no finite triangulation of the torus gives two. (b) For a finite graph with no loops or multiple edges, regarded as a one-dimensional simplicial complex, show $b_0-b_1=V-E$, including the disconnected and empty cases. Compare the connected case with Theorem~\ref{thm:wedge}.
\end{exercise}

\begin{exercise}[$\star\star$]\label{ex:coveringchi}
\emph{Counting lifted simplices.} Let $p:\tilde X\to|K|$ be a finite $k$-sheeted covering of a finite simplicial complex. (a) Show that each simplex has $k$ lifts and that their images form a triangulation of $\tilde X$. (b) Deduce $\chi(\tilde X)=k\chi(K)$. (c) Recover the index formula $1+k(n-1)$ for a finite-index subgroup of $F_n$ and the formula $2-2h=k(2-2g)$ for a connected finite cover $\Sigma_h\to\Sigma_g$.
\end{exercise}

\begin{exercise}[$\star\star$]\label{ex:lusternik}
\emph{A closed cover of a sphere.} If $n+1$ closed sets $F_1,\ldots,F_{n+1}$ cover $S^n$, prove that one contains an antipodal pair. Use distances to the first $n$ sets, allowing some sets to be empty. Include $n=0$.
\end{exercise}

\begin{exercise}[$\star$]\label{ex:pancakes}
\emph{From area to points.} (a) Two bounded measurable pancakes lie in a plane. Prove that one line bisects both areas. (b) For two finite point sets $P,Q\subset\R^2$, find a line with at most half of each set on either open side. Points on the line are allowed. Justify the limit from small equal-radius discs and include empty or singleton sets.
\end{exercise}

\begin{exercise}[$\star\star$]\label{ex:v9-degree-attachment}
\emph{A cycle whose multiple bounds.} For $d\in\Z$, attach a disc to $S^1$ by $z\mapsto z^d$, obtaining $M_d$. (a) Compute all integer homology groups of $M_d$, treating $d=0$ and negative $d$. (b) Compute them with coefficients $\Q$ and $\Z/2$. Explain how torsion invisible over $\Q$ can create both a one-dimensional and a two-dimensional class modulo two. (c) Compare the cases $d=0,1,2$ with the space, disc or projective plane described by the attachment.
\end{exercise}

\begin{exercise}[$\star\star$]\label{ex:v9-local-homology}
\emph{Recognizing a boundary or branch point.} (a) For $n\ge1$, compute $H_k(\R^n,\R^n\setminus\{0\})$ and $H_k(H,H\setminus\{0\})$, where $H=\{x_n\ge0\}$. (b) For a vertex $v$ with $d\ge1$ incident edge ends in a unit-edge graph, compute its local groups $H_k(X,X\setminus\{v\})$. Treat an isolated vertex too. (c) Deduce that a graph with a valence-three vertex is not a one-dimensional manifold with or without boundary, and that a boundary point of a manifold cannot have an interior neighbourhood.
\end{exercise}

\begin{exercise}[$\star\star$]\label{ex:v9-odd-degree}
\emph{Odd maps have odd degree.} For $n\ge1$, prove that every odd self-map $f:S^n\to S^n$ has odd integer degree. Use the transfer sequence when $n\ge2$ and lifted arguments when $n=1$. Deduce that an odd self-map is onto and cannot extend to a map $\bar B^{n+1}\to S^n$.
\end{exercise}

\begin{exercise}[$\star\star$]\label{ex:v9-torsion-matrix}
\emph{The same chains over different fields.} Consider $C_2=\Z^2$, $C_1=\Z^3$, $C_0=\Z$, with $\partial_1=0$ and
\[
\partial_2=\begin{pmatrix}4&6\\6&12\\0&0\end{pmatrix},
\]
and all other groups zero. (a) Give explicit integer row and column operations computing its homology. (b) Reduce the matrices over $\Q,\Z/2,\Z/3$ and compute the resulting homology in each degree. (c) Check the Euler characteristic in every case and explain why ranks over different fields detect information that the Euler characteristic loses.
\end{exercise}

\section{Hints}

\begin{hint}{ex:kleinrp2}
Remove the attached disc's centre and use Mayer--Vietoris as for $\Sigma_g$. The boundary circle maps to twice $a$ or twice $b$ in the bouquet's $H_1$.
\end{hint}
\begin{hint}{ex:hairydegree}
Orthogonality keeps the displayed homotopy on the sphere. For (c), normalize the straight segments from $f(x)$ to $-x$ or $x$, and check exactly when they could pass through zero.
\end{hint}
\begin{hint}{ex:lhuilier}
Count one bottom edge, one left edge and one diagonal per small square of the torus grid. For a graph use the alternating-rank formula, with no chains above dimension one.
\end{hint}
\begin{hint}{ex:coveringchi}
Lift a simplex's inclusion from a chosen point over one vertex. For two different simplices, examine lifts over their common face, not merely over an arbitrary intersection point.
\end{hint}
\begin{hint}{ex:lusternik}
Use $f_i(x)=d(x,F_i)$ when $F_i\ne\emptyset$, and $f_i(x)=1$ when it is empty. A zero distance detects membership because the set is closed.
\end{hint}
\begin{hint}{ex:pancakes}
A bisecting line leaves at most half of a family's discs wholly on one side. Pass to a convergent subsequence of normalized line equations; any point strictly on one side of the limit has its whole tiny disc on that side eventually.
\end{hint}
\begin{hint}{ex:v9-degree-attachment}
In the reduced Mayer--Vietoris sequence, the boundary map $H_1(S^1;R)\to H_1(S^1;R)$ is multiplication by $d$. Compute its kernel and cokernel over each coefficient ring.
\end{hint}
\begin{hint}{ex:v9-local-homology}
Use a small ball or star and excision. The punctured half-space retracts onto a closed hemisphere, and a punctured $d$-arm star has $d$ contractible components.
\end{hint}
\begin{hint}{ex:v9-odd-degree}
Oddness induces a map of the transfer short exact sequences. Start at $H_0$ and use naturality of the connecting maps up to the top transfer. On $S^1$, the lift difference for a half turn lies in $1/2+\Z$.
\end{hint}
\begin{hint}{ex:v9-torsion-matrix}
In the upper $2\times2$ block, subtract the first row from the second, swap them, then clear the first column and row. Over a field only the rank of this boundary map remains.
\end{hint}

\section{Solutions}

\begin{solution}{ex:kleinrp2}
\idea Integer signs in the attaching word survive as torsion. Modulo two, a twice-traversed boundary becomes zero.
\writeup (a) The punctured attached disc retracts onto the circle skeleton. As in Example~\ref{ex:surfacesH}, reduced Mayer--Vietoris gives
\[
0\to H_2(\R P^2)\to\Z\xrightarrow{\times2}\Z\to H_1(\R P^2)\to0.
\]
The doubling map is injective with cokernel $\Z/2$, so $H_2=0$ and $H_1=\Z/2$. Over $\Z/2$ the map vanishes, giving $H_2=\Z/2$ (and $H_1=\Z/2$). The space is connected, so $H_0=\Z$ over the integers.

(b) In the Klein bottle the corners identify to one vertex and the two edges form a bouquet. The boundary word has additive image $a+b-a+b=2b$. Thus
\[
0\to H_2\to\Z\xrightarrow{1\mapsto(0,2)}\Z^2\to H_1\to0.
\]
The map is injective with cokernel $\Z\oplus\Z/2$. There is no higher homology, by the same Mayer--Vietoris sequence, so the integral Betti numbers are $1,1,0$ and $\chi=0$. The torsion factor contributes no integral rank.
\end{solution}

\begin{solution}{ex:hairydegree}
\idea A nonzero tangent direction gives a continuous half-turn from each point to its antipode. Degree prevents those half-turns from existing in even dimensions.
\writeup (a) Tangency gives $x\perp X(x)$, so the displayed vector has squared norm $\cos^2(\pi t)+\sin^2(\pi t)=1$. It is continuous, equals $x$ at $t=0$ and $-x$ at $t=1$. (b) Homotopy invariance would give $1=\deg(-\mathrm{id})=(-1)^{n+1}$, forcing $n$ odd.

(c) If $f(x)\ne x$, the segment $(1-t)f(x)-tx$ never vanishes: vanishing for $0<t<1$ would make the two unit vectors $f(x),x$ equal, while endpoints cannot vanish. Its normalization therefore joins $f$ to $-\mathrm{id}$. If $f(x)\ne-x$, normalize $(1-t)f(x)+tx$ to join $f$ to the identity. For even $n$, the respective degrees would be $-1$ and $1$, so both avoidance conditions cannot hold together. On $S^0$ every tangent vector is zero, so there is no nowhere-zero tangent field; for any self-map, the image of either chosen point is that point or its antipode, giving the final alternative directly.
\end{solution}

\begin{solution}{ex:lhuilier}
\idea The torus's zero count is $1-2+1$, the alternating homology ranks. Changing the triangles cannot change those groups.
\writeup (a) Each of the nine small squares contributes its bottom, left and diagonal edges, and two triangles, after opposite-side gluing. Thus there are nine vertices, $27$ edges and $18$ triangles, giving zero. The comparison theorem makes this the same for every finite triangulation, and the surface homology calculation gives $1-2+1=0$ independently. A toroidal frame boundary therefore does not meet the sphere-boundary hypothesis in the polyhedron formula.

(b) For a one-dimensional simplicial complex, the only simplex counts are $V,E$ and the only possible Betti numbers are $b_0,b_1$. Euler--Poincar\'e gives $b_0-b_1=V-E$. For a nonempty connected graph $b_0=1$, so $b_1=E-V+1$, matching the free-group rank in Theorem~\ref{thm:wedge}. For a disconnected graph the ranks add over components. For the empty graph all four quantities are zero.
\end{solution}

\begin{solution}{ex:coveringchi}
\idea Lift each whole simplex, then use uniqueness on its faces to glue the lifted pieces. This multiplies every simplex count by the number of sheets.
\writeup (a) A simplex is convex and locally path-connected, so its inclusion into $|K|$ has a unique lift through any chosen point over one vertex, by the lifting criterion. There are $k$ choices. Each lift is injective because projecting it recovers the original simplex. If two lifts of that simplex meet, their meeting point projects to the same domain point, so uniqueness makes them identical. The restrictions of these lifts to a face are precisely the lifts of that face.

For two downstairs simplices $\sigma,\eta$, their intersection is their common face or empty. Over a nonempty common face, two chosen lifts either coincide on that whole face or have disjoint images, by uniqueness. Thus two lifted simplex images intersect in a common lifted face or not at all. Their vertex sets also distinguish them: equal lifted vertex sets project to the same simplex and give agreement at a vertex. They form a finite simplicial complex $\tilde K$, ordered by the downstairs order with arbitrary tie breaking within fibres.

The glued lift map $|\tilde K|\to\tilde X$ is a continuous bijection. The total space is Hausdorff: different projections have disjoint inverse-image neighbourhoods, and points over the same basepoint separate into different sheets of an evenly covered neighbourhood. Compactness of the finite realization now makes the bijection a homeomorphism. If $k=0$ is allowed, the total space and lifted complex are empty and the same conclusion holds.

(b) Each $c_j$ multiplies by $k$, hence so does its finite alternating sum. (c) Subdivide each circle in $R_n$ into three edges to get an actual simplicial complex with $\chi=1-n$. Its connected $k$-sheeted graph cover has $\chi=1-\operatorname{rank}H$, giving the index formula. For closed orientable surfaces, the quoted triangulation theorem and $\chi(\Sigma_g)=2-2g$ give $2-2h=k(2-2g)$. These connected nonempty covers have $k\ge1$.
\end{solution}

\begin{solution}{ex:lusternik}
\idea Equal distances at antipodes force either a shared membership in one of the first sets or membership in the last set for both points.
\writeup For $n\ge1$, put $f_i(x)=d(x,F_i)$ for nonempty $F_i$, and $f_i(x)=1$ otherwise, for $1\le i\le n$. Each coordinate is continuous. Borsuk--Ulam supplies $x$ with $f(x)=f(-x)$. If some coordinate is zero, that set is nonempty and closed, so both points belong to it. If every coordinate is positive, neither point belongs to any of the first $n$ sets. The covering assumption puts both in $F_{n+1}$. For $n=0$, the only covering set is all of $S^0$, which already contains the antipodal pair.
\end{solution}

\begin{solution}{ex:pancakes}
\idea Area bisections of shrinking discs give a limiting line. Points strictly on either side of that line eventually contribute whole discs there, which bounds their number.
\writeup (a) Apply the ham sandwich theorem with $n=2$. (b) If $P\cup Q$ is empty, any line works. If it is a singleton, any line through that point works, whether the point belongs to one set or both. Otherwise choose $\varepsilon_j\downarrow0$ smaller than half the minimum separation of distinct points in $P\cup Q$. For each family replace its points by discs of radius $\varepsilon_j$. Within each family these are disjoint equal-area discs. A point common to the two families simply supplies a disc to each.

Part (a) gives a line $\ell_j=\{x:\langle u_j,x\rangle=c_j\}$, $|u_j|=1$, bisecting both families' areas. At least one family is nonempty; a line bisecting its positive area meets a fixed bounded disc containing all the small discs, so $c_j$ is bounded. Pass to a subsequence with $(u_j,c_j)\to(u,c)$, where $|u|=1$.

Every point strictly on a given side of $\ell=\{\langle u,x\rangle=c\}$ eventually has its whole tiny disc on that side of $\ell_j$. Since there are finitely many points, this holds simultaneously for all points of either family strictly on that side. If there are $s$ of them in $P$, their disjoint discs contribute $s\pi\varepsilon_j^2$ to a half-area of $|P|\pi\varepsilon_j^2/2$, so $s\le|P|/2$. The same argument applies to $Q$ and the other side. Empty families cause no restriction. Points on the limiting line need not be counted on either open side. The cited algorithmic work uses this discrete bisection principle.
\end{solution}

\begin{solution}{ex:v9-degree-attachment}
\idea The attached face kills $d$ times the circle generator. Its kernel measures top homology; its cokernel measures which circle classes remain.
\writeup (a) Cover $M_d$ by the complement $U$ of the disc's centre and the open disc $V$. The first retracts onto the original circle, the second contracts, and their intersection retracts onto a circle. Its generator maps to the attaching loop of degree $d$, so Mayer--Vietoris gives
\[
0\to H_2(M_d;R)\to R\xrightarrow{\times d}R\to H_1(M_d;R)\to0.
\]
The space is connected, so $H_0(M_d;R)=R$, and the same sequence makes all groups above degree two zero. Over $\Z$, if $d\ne0$ the kernel is zero and the cokernel is $\Z/|d|$ (zero for $|d|=1$). If $d=0$, both $H_1$ and $H_2$ are $\Z$.

(b) Over $\Q$, nonzero $d$ is invertible, giving $H_1=H_2=0$; for $d=0$ both are $\Q$. Over $\Z/2$, odd $d$ gives an invertible map and zero groups, while even $d$, including zero and negative even integers, gives the zero map and $H_1=H_2=\Z/2$. Thus a nonzero even attaching degree has no positive-dimensional rational homology, but modulo two both classes appear. The equation that used to identify a multiple of the loop with a boundary becomes zero, leaving both the loop and the attached face as classes.

(c) For $d=0$ the attaching circle collapses to one point of the original circle. The quotient disc with its whole boundary collapsed is $S^2$, so $M_0=S^1\vee S^2$, consistent with the two free classes. For $d=1$ the boundary is attached homeomorphically to the original circle, producing a disc and zero reduced homology. For $d=2$, each pair of antipodal boundary points has the same image; the circle is identified with that boundary quotient, so $M_2$ is the projective plane. Its order-two integer class and mod-two top class agree with the earlier calculation.
\end{solution}

\begin{solution}{ex:v9-local-homology}
\idea Local homology records how many cycles remain when one point is removed from a small neighbourhood. A boundary point and an interior point leave different punctured neighbourhoods.
\writeup (a) Both $\R^n$ and $H$ are contractible. Their punctured versions are nonempty for $n\ge1$, so the pair sequences give zero in degree zero and identify relative degree $k\ge1$ with reduced degree $k-1$ of the punctured space. Radial projection retracts $\R^n\setminus\{0\}$ onto $S^{n-1}$, giving
\[
H_k(\R^n,\R^n\setminus\{0\})=
\begin{cases}\Z,&k=n,\\0,&k\ne n.\end{cases}
\]
For $H\setminus\{0\}$ the radial image is the closed upper hemisphere. It contracts to its north pole, by normalized straight segments, including the single-point hemisphere when $n=1$. Every relative group for $(H,H\setminus\{0\})$ is therefore zero.

(b) Excise the complement of a small open star $S$ about $v$. That complement is closed and contained in the open subspace $X\setminus\{v\}$, so excision applies. The star contracts to $v$, while $S\setminus\{v\}$ has $d$ contractible interval components. Its reduced $H_0$ is $\Z^{d-1}$ and all other reduced groups vanish. The pair sequence gives local $H_1=\Z^{d-1}$ and zero in every other degree. For an isolated vertex the local pair is $(\{v\},\emptyset)$, giving $H_0=\Z$ and zero otherwise.

(c) In a one-manifold an interior point has local $H_1=\Z$, and a boundary point has all local groups zero by (a). Valence three gives $\Z^2$, so no such manifold chart exists there. More generally, in positive dimension an interior point has nonzero top local homology and a boundary point has none. A homeomorphism of neighbourhoods preserves these pair groups by excision and functoriality, so it cannot interchange the two types. Dimension zero has only isolated interior points and no boundary charts of positive half-space type.
\end{solution}

\begin{solution}{ex:v9-odd-degree}
\idea The transfer sequence forces an odd map to preserve the top class modulo two. An integer multiple preserving that class must be odd.
\writeup For $n\ge2$, oddness induces $\bar f:\R P^n\to\R P^n$ and a map of the transfer short exact sequences, since $f$ carries antipodal lifts to antipodal lifts. On $H_0(\R P^n;\Z/2)$ its action is the identity. Naturality and the connecting isomorphisms proved in Borsuk--Ulam make it the identity in every degree through $n$. The top transfer is an isomorphism, so $f_*$ on $H_n(S^n;\Z/2)$ is the identity as well.

To relate this to integer degree, reduce coefficients modulo two. The natural map $H_n(S^n;\Z)\to H_n(S^n;\Z/2)$ sends a generator to the nonzero generator: this follows inductively from the natural sphere connecting isomorphisms, starting with $[1]-[-1]$ on $S^0$. If the integer action is multiplication by $D$, its mod-two action is therefore multiplication by $D\bmod2$. Hence $D$ is odd.

For $n=1$, write $f(e^{2\pi it})=e^{2\pi i\alpha(t)}$ with a continuous lift on $[0,1]$. Oddness gives $\alpha(t+1/2)-\alpha(t)\in1/2+\Z$ for $0\le t\le1/2$. This difference is constant, so the total change is an odd integer. It is the winding number and hence the degree.

An odd degree is nonzero, so $f$ cannot miss a point. If $f$ extended to $F:\bar B^{n+1}\to S^n$, its top reduced homology action would factor as $f_*=F_*i_*$ through the zero reduced homology of the ball, forcing degree zero. Thus no extension exists.
\end{solution}

\begin{solution}{ex:v9-torsion-matrix}
\idea The integer diagonal entries retain finite-order classes. Over a field each diagonal entry either becomes invertible or vanishes, changing both kernel and cokernel.
\writeup (a) All differentials compose to zero because $\partial_1=0$. In the upper block, perform $R_2\leftarrow R_2-R_1$, interchange $R_1,R_2$, then $R_2\leftarrow R_2-2R_1$. This gives
\[
\begin{pmatrix}4&6\\6&12\end{pmatrix}
\longrightarrow\begin{pmatrix}4&6\\2&6\end{pmatrix}
\longrightarrow\begin{pmatrix}2&6\\4&6\end{pmatrix}
\longrightarrow\begin{pmatrix}2&6\\0&-6\end{pmatrix}.
\]
Next use $C_2\leftarrow C_2-3C_1$ and change the second row's sign to obtain $\operatorname{diag}(2,6)$, with the third row still zero. These are invertible integer basis changes. Thus $H_2=0$, $H_1=\Z\oplus\Z/2\oplus\Z/6$, and $H_0=\Z$. All other groups vanish.

(b) The same operations remain invertible over each indicated field. Over $\Q$, the boundary has rank two, giving $(H_0,H_1,H_2)=(\Q,\Q,0)$. Modulo two it is zero, giving $((\Z/2),(\Z/2)^3,(\Z/2)^2)$. Modulo three its diagonal ranks reduce to one, giving $((\Z/3),(\Z/3)^2,\Z/3)$. In each case these also follow directly from $H_1=\operatorname{coker}\partial_2$ and $H_2=\ker\partial_2$.

(c) The chain count is $1-3+2=0$. The integral and rational alternating ranks are $1-1=0$, the mod-two sum is $1-3+2=0$, and the mod-three sum is $1-2+1=0$. Changes in a boundary rank add the same amount to adjacent homology dimensions, so they cancel in this alternating sum. The groups reveal two independent mod-two kernel classes but only one mod-three class, information that a single Euler characteristic cannot distinguish.
\end{solution}

\section*{Sources and further reading}
\addcontentsline{toc}{section}{Sources and further reading}

Euler's letter to Goldbach of 14 November 1750 is in \emph{Leonhard Euler und Christian Goldbach: Briefwechsel 1729--1764}, ed.\ A.~P.~Juškevič and E.~Winter (Akademie-Verlag, 1965); his proofs
are ``Elementa doctrinae solidorum'' and ``Demonstratio nonnullarum insignium proprietatum quibus solida hedris planis inclusa sunt praedita'', \emph{Novi Comm.\ Acad.\ Sci.\ Petrop.} 4 (1752/53, printed 1758),
109--140 and 140--160. S.~Lhuilier, ``M\'emoire sur la poly\'edrom\'etrie'', \emph{Ann.\ Math.\ Pures Appl.} 3 (1812--13), 169--189. I.~Lakatos, \emph{Proofs and Refutations} (Cambridge University Press, 1976);
D.~Richeson, \emph{Euler's Gem} (Princeton University Press, 2008), tells the whole story. On Emmy Noether: P.~Alexandroff and H.~Hopf, \emph{Topologie I} (Springer, 1935), preface; H.~Hopf, ``Eine
Verallgemeinerung der Euler--Poincar\'eschen Formel'', \emph{Nachr.\ Ges.\ Wiss.\ G\"ottingen} (1928), 127--136; C.~McLarty, ``Emmy Noether's set-theoretic topology'', in \emph{The Architecture of Modern
Mathematics} (Oxford University Press, 2006).

W.~Mayer, ``\"Uber abstrakte Topologie'', \emph{Monatsh.\ Math.\ Phys.} 36 (1929), 1--42 and 219--258; L.~Vietoris, ``\"Uber die Homologiegruppen der Vereinigung zweier Komplexe'', \emph{Monatsh.\ Math.\ Phys.}
37 (1930), 159--162. L.~E.~J.~Brouwer, ``Beweis der Invarianz des $n$-dimensionalen Gebiets'', \emph{Math.\ Ann.} 71 (1912), 305--313. K.~Borsuk, ``Drei S\"atze \"uber die $n$-dimensionale euklidische Sph\"are'',
\emph{Fund.\ Math.} 20 (1933), 177--190; A.~H.~Stone and J.~W.~Tukey, ``Generalized `sandwich' theorems'', \emph{Duke Math.\ J.} 9 (1942), 356--359; W.~A.~Beyer and A.~Zardecki, ``The early history of the ham sandwich theorem'', \emph{Amer.\ Math.\ Monthly} 111 (2004), 58--61, which translates the 1938 note. C.-Y.~Lo, J.~Matou\v sek and W.~Steiger, ``Algorithms for ham-sandwich cuts'', \emph{Discrete Comput.\ Geom.} 11 (1994), 433--452.

A.~Hatcher, \emph{Algebraic Topology} (Cambridge University Press, 2002), Sections~2.1--2.2 and~2.B, gives further details on small chains, invariance of domain and the transfer argument. The simplex-by-simplex comparison here lets you return to finite matrices after learning the singular theory. For concrete antipodal and partition problems, continue with J.~Matou\v sek, \emph{Using the Borsuk--Ulam Theorem} (Springer, 2003).
